%% supplementary_material.tex
%% Online Supplementary Material
%% Contains material moved from the main text during R2 revision for word count compliance.

\section*{Online Supplementary Material}

This supplementary document contains background and methodological detail moved from the main text to meet journal word count requirements. Section and paragraph headings correspond to the main manuscript locations from which the material was relocated.

\subsection*{From Section 3.1: Coal Price Cycle and Export Markets}

\textit{Recent coal price cycle.} The recent coal price cycle provides important context for understanding the sector's transition risk exposure. Newcastle benchmark prices collapsed to approximately \$48/tonne in September 2020 as COVID-19 crushed energy demand, then surged to unprecedented levels above \$400/tonne in mid-2022 following Russia's invasion of Ukraine and the ensuing global energy security panic. By 2023--2025, prices normalised to the \$100--140/tonne range as supply adjusted and renewable energy deployment accelerated. This extreme volatility illustrates the commodity price risk to which Indonesian coal producers and their bank lenders are exposed, while the long-run downward trend in prices is broadly consistent with the transition pathways modelled in this paper. Notably, China's unofficial ban on Australian coal imports from late 2020 temporarily redirected demand toward Indonesian suppliers, reinforcing Indonesia's role as the swing exporter in the seaborne thermal coal market.

\textit{Export markets and competitive position.} Indonesian thermal coal is exported primarily to China, India, Japan, South Korea, and Southeast Asian neighbours. Its competitive advantage rests on geographic proximity to Asian demand centres and low extraction costs, but it faces growing competition from Australian and South African coal at the higher-calorie end and from Russian coal that has been redirected to Asian markets following Western sanctions. Under aggressive decarbonisation scenarios, demand from Japan and South Korea, both of which have announced net-zero targets, is projected to decline sharply by 2035, while Chinese and Indian demand paths remain subject to considerable uncertainty \citep{IEA2023, Mendelevitch2019}.

\subsection*{From Section 3.3: Sub-National Fiscal Detail}

Beyond the central government, coal-producing provinces and regencies in Kalimantan and South Sumatra receive revenue-sharing transfers (Dana Bagi Hasil) from coal royalties, which constitute a major source of sub-national government funding. The key to the government revenue is dividends from State-owned coal enterprise. PT Bukit Asam (PTBA), as the only publicly listed state-owned coal company, pays dividends to the government through the Ministry of State-Owned Enterprises. With a market capitalisation of approximately \$1.9 billion and a historical dividend yield of roughly 7\%, PTBA contributes an estimated \$130 million annually in dividend income to the state. While this is modest relative to total revenues, it represents a direct and visible fiscal stake in coal sector profitability that complicates political decisions regarding coal phase-down.

\subsection*{From Section 3.4: Regulatory Landscape Detail}

\textit{JETP implementation detail.} The Just Energy Transition Partnership concept was first announced at COP26 in Glasgow (November 2021), with Indonesia's partnership formally launched at COP27 in Sharm el-Sheikh (November 2022). The \$20 billion financing package, comprising \$10 billion from the International Partners Group (G7 plus Norway) and a matched private finance mobilisation target, is channelled through an Energy Transition Mechanism (ETM) designed to accelerate early coal plant retirement and scale up renewable energy capacity. The ETM pipeline, managed jointly by the Asian Development Bank and the Indonesian government, has identified specific coal plants for early retirement and blended finance facilities. However, as of early 2026, implementation has been slower than anticipated: the financing arrangements remain under negotiation, just-transition provisions for affected workers and communities are not fully specified, and the role of Indonesian banks in refinancing retired coal assets has yet to be formalised.

\textit{Bank Indonesia green taxonomy detail.} Bank Indonesia (BI) has taken complementary steps at the monetary policy level. The Bank Indonesia Sustainable Finance Taxonomy (Green Taxonomy), published in 2022 and aligned with ASEAN Taxonomy for Sustainable Finance, classifies economic activities by their environmental impact and provides a framework for green bond eligibility and preferential refinancing treatment. BI's 2023 circular (SEBI No. 14/2023) encourages banks to disclose their exposure to taxonomy-classified activities, creating an indirect incentive to reduce brown lending. BI has also piloted a green repo facility, offering preferential rates on repurchase agreements collateralised by green bonds, though the scale remains limited relative to the banking system's total assets.

\subsection*{From Section 6: NPL Shock Derivation Detail}

The NPL shock calibration warrants discussion. The 15 percentage point increment under the 1.5\textdegree C scenario is applied to the \textit{residual} mining portfolio, which includes non-coal mining (nickel, tin, gold) that may not be directly affected by the coal transition. This is a conservative (loss-increasing) assumption. For context, the historical peak NPL ratio in the Indonesian mining sector was 7.7\% during the 2015--2016 commodity price downturn, when Newcastle coal prices fell to approximately \$48/tonne. Our 1.5\textdegree C NPL shock of 15 percentage points thus represents roughly twice the worst historical episode, reflecting the deeper and more sustained nature of the structural transition relative to a cyclical price downturn.

\subsection*{From Section 8.3: Comparative Country Detail}

\textit{India.} India represents the largest single country coal exposure globally, with approximately 60\% of electricity from coal and a rapidly growing economy whose energy demand projections create acute policy uncertainty. Indian public sector banks have substantially larger coal-related exposures than their Indonesian counterparts in absolute terms, and the Indian banking system's pre-existing capital quality challenges (elevated NPLs, recapitalisation requirements) mean that transition-related credit losses would arrive on already-stressed balance sheets. However, India's larger and more diversified economy implies smaller GDP-level fiscal impacts per unit of coal revenue loss.

\textit{Colombia and the Philippines.} Colombia and the Philippines represent mid-tier cases where coal export dependence creates fiscal vulnerability without the large domestic financial system that makes Indonesia and South Africa particularly interesting from a banking stability perspective. In Colombia, the coal sector's contribution to export earnings (approximately 12\%) and government revenues is comparable to Indonesia's, but Colombia's financial system is smaller relative to GDP and the bank-coal lending relationship is less concentrated in domestic institutions. The Philippines has a smaller domestic coal sector but significant coal-fired power exposure, creating transition risk through the utilities sector rather than the mining sector directly.
